\chapter{Possibili miglioramenti}\label{chapter:migliorie}
\section{Lavori in corso}\label{sec:cap_sec_subsec_lavori_in_corso}

La ristrutturazione del gestionale è attualmente in fase di sviluppo, poiché non è ancora completamente finalizzato. Nonostante le funzionalità principali siano state implementate e siano operative, ci sono ancora diverse sezioni da completare, ridefinire o modificare. Questi aggiornamenti sono essenziali per migliorare l'efficienza del sistema e per garantire che risponda pienamente alle esigenze aziendali di Dieffetech.

Il processo di revisione continua assicura che il rifacimento del gestionale si evolva in modo dinamico, migliorando costantemente la sua usabilità e la sua capacità di supportare il team. Questo approccio iterativo ci permette di affinare ogni dettaglio e di rilasciare nuove versioni con funzionalità aggiuntive, aumentando gradualmente il valore complessivo del sistema. 

\section{AI}\label{sec:cap_sec_subsec_AI}

Un potenziale miglioramento del gestionale potrebbe essere rappresentato dall’integrazione di strumenti basati sull’Intelligenza Artificiale (AI), in grado di ottimizzare i processi aziendali. In particolare, l'adozione di algoritmi di analisi predittiva consentirebbe di analizzare i dati storici relativi a clienti e lead, fornendo suggerimenti utili ai dipendenti per l'adozione di decisioni più informate. Tali algoritmi potrebbero, ad esempio, identificare con maggiore precisione i lead con maggiore probabilità di conversione o le aree aziendali che richiedono interventi strategici.

In aggiunta, l'AI potrebbe essere impiegata per l'automazione di attività ripetitive, come la generazione automatica di report periodici o la classificazione e l’archiviazione di documenti e allegati. Queste soluzioni ridurrebbero il carico di lavoro manuale dei dipendenti, permettendo loro di concentrarsi su attività a più alto valore aggiunto. Un ulteriore sviluppo potrebbe includere l'integrazione di assistenti virtuali per supportare il personale nell'utilizzo del gestionale, semplificando e velocizzando l'accesso alle informazioni e la gestione dei processi aziendali.

\section{Aggiornamenti e aggiunte}\label{sec:cap_sec_subsec_aggiornamenti_e_aggiunte}

Uno dei punti di forza di un gestionale interno risiede nella sua capacità di adattarsi alle nuove esigenze aziendali che si presentano nel tempo. Il sistema, infatti, è progettato per essere facilmente aggiornabile e ampliabile in base alle necessità operative e strategiche dell'azienda. L'approccio modulare consente l'integrazione di nuove funzionalità, il miglioramento delle prestazioni e la risoluzione di bug, senza interrompere il flusso di lavoro.

In un contesto in continua evoluzione come quello di Dieffetech, il gestionale è soggetto a costanti aggiornamenti per riflettere cambiamenti nelle procedure aziendali, l'introduzione di nuovi strumenti tecnologici o l'integrazione di soluzioni innovative. Inoltre, aggiunte future potrebbero riguardare nuove sezioni o moduli specifici, pensati per ottimizzare ulteriormente la gestione di clienti o altri processi aziendali.

Il continuo sviluppo del gestionale rappresenta un elemento chiave per garantire che il sistema resti in linea con gli obiettivi aziendali e che supporti efficacemente i dipendenti nelle loro attività quotidiane. Questo assicura che il gestionale mantenga la sua flessibilità e scalabilità nel tempo, fornendo soluzioni sempre più aderenti alle esigenze operative dell'azienda.